\documentclass[runningheads]{llncs}
% \usepackage[pagebackref]{hyperref}

% ---------------------------------------------------------------
% Include basic ECCV package
 
% TODO REVIEW: Insert your submission number below by replacing '*****'
% TODO FINAL: Comment out the following line for the camera-ready version
\usepackage[review,year=2026,ID=1259]{eccv}
% TODO FINAL: Un-comment the following line for the camera-ready version
%\usepackage{eccv}

% OPTIONAL: Un-comment the following line for a version which is easier to read
% on small portrait-orientation screens (e.g., mobile phones, or beside other windows)
%\usepackage[mobile]{eccv}


% ---------------------------------------------------------------
% Other packages

% Commonly used abbreviations (\eg, \ie, \etc, \cf, \etal, etc.)
\usepackage{eccvabbrv}
% \usepackage[pagebackref]{hyperref}

% Include other packages here, before hyperref.
\usepackage{graphicx}
\usepackage{booktabs}

% The "axessiblity" package can be found at: https://ctan.org/pkg/axessibility?lang=en
\usepackage[accsupp]{axessibility}  % Improves PDF readability for those with disabilities.
\definecolor{myblue}{RGB}{193,248,255}

% ---------------------------------------------------------------
% Hyperref package

% It is strongly recommended to use hyperref, especially for the review version.
% Please disable hyperref *only* if you encounter grave issues.
% hyperref with option pagebackref eases the reviewers' job, but should be disabled for the final version.
%
% If you comment hyperref and then uncomment it, you should delete
% main.aux before re-running LaTeX.
% (Or just hit 'q' on the first LaTeX run, let it finish, and you
%  should be clear).

% TODO FINAL: Comment out the following line for the camera-ready version
\usepackage[pagebackref,breaklinks,colorlinks,citecolor=eccvblue]{hyperref}
% TODO FINAL: Un-comment the following line for the camera-ready version
% \usepackage{hyperref}
\usepackage{booktabs,multirow}
\usepackage{makecell}
\usepackage{array}
\usepackage[table]{xcolor}
\usepackage{xspace}
\usepackage{graphicx}
\usepackage{amsmath}
\usepackage{amssymb}
\usepackage{caption}
%\usepackage[compatibility=false]{caption}
\usepackage{subcaption}
% Support for ORCID icon
\usepackage{orcidlink}
\definecolor{colorNormal}{RGB}{193,248,255} % 浅蓝色 (Normal 最佳，与myblue一致)
\definecolor{colorDisco}{RGB}{255,235,205}  % 浅橙色 (Disco 最佳)
\definecolor{colorTwo}{RGB}{193,248,255} % 浅蓝色 (Single View 最佳)
\definecolor{colorSingle}{RGB}{255,235,205}  % 浅橙色 (Two Views 最佳)

\begin{document}

% ---------------------------------------------------------------
% TODO REVIEW: Replace with your title
\title{Supplementary Materials for R3D: Revisiting 3D Policy Learning} 

% TODO REVIEW: If the paper title is too long for the running head, you can set
% an abbreviated paper title here. If not, comment out.
\titlerunning{Abbreviated paper title}

% TODO FINAL: Replace with your author list. 
% Include the authors' OCRID for the camera-ready version, if at all possible.
\author{First Author\inst{1}\orcidlink{0000-1111-2222-3333} \and
Second Author\inst{2,3}\orcidlink{1111-2222-3333-4444} \and
Third Author\inst{3}\orcidlink{2222--3333-4444-5555}}

% TODO FINAL: Replace with an abbreviated list of authors.
\authorrunning{F.~Author et al.}
% First names are abbreviated in the running head.
% If there are more than two authors, 'et al.' is used.

% TODO FINAL: Replace with your institution list.
\institute{Princeton University, Princeton NJ 08544, USA \and
Springer Heidelberg, Tiergartenstr.~17, 69121 Heidelberg, Germany
\email{lncs@springer.com}\\
\url{http://www.springer.com/gp/computer-science/lncs} \and
ABC Institute, Rupert-Karls-University Heidelberg, Heidelberg, Germany\\
\email{\{abc,lncs\}@uni-heidelberg.de}}

\maketitle
\begin{table}[]
\vspace{-10pt}
    \centering
    \small
    \renewcommand{\arraystretch}{1.1} % 紧凑行高
    \setlength{\tabcolsep}{3.2pt}      % 进一步压缩列间距
    \caption{\textbf{Success Rates (\%) on 6 Real-World Tasks in 4 Types.}}
    % \vspace{-1em}
    \label{tab:realworld_structured}
    
    \begin{tabular}{l@{\hskip 3pt}cccccc|c} % 加大第一列后的间距，区分模型和数据
        \toprule
        & \textbf{Long-Horizon} & \multicolumn{3}{c}{\textbf{Pick \& Place}} & \textbf{Articulated} & \textbf{Soft Body} & \\
        \cmidrule(lr){2-2} \cmidrule(lr){3-5} \cmidrule(lr){6-6} \cmidrule(lr){7-7}
        Model & \begin{tabular}[c]{@{}c@{}}Stack Three\\Blocks\end{tabular} & \begin{tabular}[c]{@{}c@{}}Place\\Cup\end{tabular} & \begin{tabular}[c]{@{}c@{}}Fit\\Banana\end{tabular} & \begin{tabular}[c]{@{}c@{}}Place\\Kettle\end{tabular} & \begin{tabular}[c]{@{}c@{}}Open\\Drawer\end{tabular} & \begin{tabular}[c]{@{}c@{}}Fold\\Towel\end{tabular} & \textbf{Avg.} \\
        \midrule
        DP3~\cite{ze20243d}  & 2  & 12 & 42 & 40 & 28 & 54 & 29.7 \\
        DP~\cite{chi2024diffusionpolicyvisuomotorpolicy}   & \cellcolor{myblue}\textbf{12} & 58 & 48 & 64 & 46 & 36 & 44.0 \\
        \textbf{Ours} & \cellcolor{myblue}\textbf{12} & \cellcolor{myblue}\textbf{82} & \cellcolor{myblue}\textbf{64} & \cellcolor{myblue}\textbf{76} & \cellcolor{myblue}\textbf{64} & \cellcolor{myblue}\textbf{66} & \cellcolor{myblue}\textbf{60.7} \\
        \bottomrule
    \end{tabular}
    \vspace{-2em}
\end{table}
 \section{More Real-world Experiments}
To further illustrate the performance of our approach across more diverse tasks, we provide supplementary experiments on three additional tasks beyond those in the main paper, and summarize the results on all 6 tasks in Tab.~\ref{tab:realworld_structured}. As demonstrated, our approach achieves consistently strong performance across a total of six tasks spanning four distinct task types.
 
\section{Real-World Implementation Details}
% # 各种details和叙述可以仿照之前cvpr补充材料的风格
\subsection{Task Details for 6 Tasks (with \textbf{4 distinct} Task Types)}
% # 真机场景的setting和之前基本一样，这里也要插入图片，区别只是在于我们目前只有right和hand camera
% 要提一下对于老的三个任务我们有三个相机，新的两个，不过我们只贴新的图片，并且新增的几个baseline本来也只是在新的三个任务上做的

% 

% # 几个任务的描述 【Task描述(init范围+成功判定)】
% # 新的三个任务对过之后的准确的描述

% #柜子：
% #成功：把柜子拉出足够远
% #典型失败：没抓住把手，把柜子只拉出很短距离
% #摆放范围：70cmx20cm的矩形区域+20cmx20cm的矩形区域组成的L字形，角度旋转90°

% #水壶：
% #成功：把水壶放在炉子上，使水壶圆心在炉子上的蓝色区域
% #典型失败：没抓住水壶，放置位置偏移
% #摆放范围：20cmx20cm的矩形区域，角度面朝前方左右各旋转60°内

% #毛巾：
% #成功：把毛巾对折
% #典型失败：没抓住毛巾，折叠后毛巾两边不对齐
% #摆放范围：桌布上70cmx50cm的区域，角度平行于桌缘，左右各旋转30°内

\paragraph{Tasks description and success criteria:}


% \begin{itemize}
%     \item \textbf{Place Kettle on Stove}: A kettle is randomly positioned within a 20\,cm $\times$ 20\,cm rectangular region, with orientation facing forward and rotating within $\pm$60$^\circ$. \textbf{Success:} the kettle is grasped by its handle, placed on the stove such that the kettle center aligns with the blue target region on the stove, and released stably. \textbf{Typical failures:} failure to grasp the kettle handle, placement position offset.
    
%     \item \textbf{Open Drawer}: A drawer stack (three drawers) is placed within an L-shaped region composed of a 70\,cm $\times$ 20\,cm rectangle plus a 20\,cm $\times$ 20\,cm rectangle, with orientation rotated 90$^\circ$. The front-facing drawer must be opened. \textbf{Success:} the drawer is grasped by its handle and pulled open to at least 50\% of its full extension. \textbf{Typical failures:} failure to grasp the handle, drawer pulled only a short distance.
    
%     \item \textbf{Fold Towel}: A towel is randomly positioned on a tablecloth within a 70\,cm $\times$ 50\,cm region, with orientation parallel to the table edge and rotating within $\pm$30$^\circ$. \textbf{Success:} the towel is grasped and folded such that both sides align. \textbf{Typical failures:} failure to grasp the towel, misaligned sides after folding.
% \end{itemize}

\begin{itemize}

    \item \textbf{Place Cup}: A vertically standing cup is initially placed in the right-half region of a 60\,cm $\times$ 40\,cm foam board, with yaw uniformly randomized over $360^\circ$ around the axis normal to the board. \textbf{Success:} the cup is grasped and placed upright into either the left or right target hole and released stably. \textbf{Typical failures:} the cup misses the hole (stops short or collides with the rim), or tips over instead of standing inside the hole.
    
    \item \textbf{Fit Banana on Plate}: A plate is fixed on the same 60\,cm $\times$ 40\,cm foam board, centered on the mid-line and shifted 10\,cm to the right; a banana-shaped object is placed 10--15\,cm to the left of the plate, with yaw uniformly randomized over $360^\circ$. \textbf{Success:} the banana is grasped, placed fully inside the plate area, and released without touching the outer boundary. \textbf{Typical failures:} the banana is not fully on the plate, drifts to the edge, or collides with the plate rim.
    
    \item \textbf{Stack Three Blocks}: Three wooden blocks are initialized near the centers of three 20\,cm $\times$ 40\,cm cells obtained by partitioning the 60\,cm $\times$ 40\,cm foam board along its long side; each block center is jittered within $\pm$5\,cm, and block orientations remain nearly fixed. \textbf{Success:} all three blocks are sequentially transported to the central stacking region to form a stable three-block tower. \textbf{Typical failures:} misaligned placement causing tower collapse, leaving one block unmoved, or stopping after only a partial stack.

    \item \textbf{Place Kettle on Stove}: A kettle is randomly positioned within a 20\,cm $\times$ 20\,cm rectangular region, with orientation facing forward and rotating within $\pm$60$^\circ$. \textbf{Success:} the kettle is grasped by its handle, placed on the stove such that the kettle center aligns with the blue target region on the stove, and released stably. \textbf{Typical failures:} failure to grasp the kettle handle, placement position offset.
    
    \item \textbf{Open Drawer}: A drawer stack (three drawers) is placed within an L-shaped region composed of a 70\,cm $\times$ 20\,cm rectangle plus a 20\,cm $\times$ 20\,cm rectangle, with orientation rotated 90$^\circ$. The front-facing drawer must be opened. \textbf{Success:} the drawer is grasped by its handle and pulled open to at least 50\% of its full extension. \textbf{Typical failures:} failure to grasp the handle, drawer pulled only a short distance.
    
    \item \textbf{Fold Towel}: A towel is randomly positioned on a tablecloth within a 70\,cm $\times$ 50\,cm region, with orientation parallel to the table edge and rotating within $\pm$30$^\circ$. \textbf{Success:} the towel is grasped and folded such that both sides align. \textbf{Typical failures:} failure to grasp the towel, misaligned sides after folding.
\end{itemize}

\paragraph{Episode Termination and Failure Definition.}
Each trial runs for up to \textbf{60 seconds}. A trial is marked as \textbf{failure} if:
\begin{itemize}
    \item the robot becomes stuck with no meaningful motion, or    
    \item the task is not successfully completed within the 60-second limit.
\end{itemize}



\paragraph{Evaluation Protocol.}
Each method is evaluated for \textbf{50 trials per task}.
All methods share identical sensing, fusion, cropping, and control settings.

\subsection{Point cloud processing and visualization}
% # 沿用之前的图片，不过里面提到的相机的名字要和正文对齐，只使用之前的right和hand的camera，删掉left的那张图
\label{sec:realworld_all}


\begin{figure*}[tb]
    \centering
    \small

    % Row 1
    \begin{tabular}{cc}
        \includegraphics[width=0.35\textwidth]{figures/rear.png} &
        \includegraphics[width=0.35\textwidth]{figures/hand.png} \\
        (a) Eye-to-hand view (base frame) & 
        (b) Eye-in-hand view (base frame)
    \end{tabular}

    \vspace{6pt}

    % Row 2
    \begin{tabular}{cc}
        \includegraphics[width=0.35\textwidth]{figures/merged.png} &
        \includegraphics[width=0.35\textwidth]{figures/merged-zoom.png} \\
        (c) Merged + cropped (8192 pts) &
        (d) Zoom-in view of training point cloud
    \end{tabular}

    \caption{\textbf{Point cloud processing pipeline.}
    (a--b) point clouds from eye-to-hand view and eye-in-hand view are individually back-projected and transformed into the base frame. (c) Merged and cropped point cloud downsampled to 8,192 points via FPS during training. (d) Zoom-in view showing geometric details.}
    \label{fig:train_pc_all}
\end{figure*}
\paragraph{Point Cloud Processing Pipeline.}
For each demonstration episode, we reconstruct a unified colored point cloud in the robot base frame using a standardized multi-view RGBD pipeline:

\begin{enumerate}
    \item \textbf{RGBD information pre-processing}: For each camera, raw RGB and (optionally CDM(\textbf{Camera Depth Model}~\cite{liu2025manipulation})-enhanced) depth maps are undistorted and back-projected into 3D using calibrated intrinsics.
    
    \item \textbf{Unified coordinate transformation}: Point clouds from eye-to-hand view and eye-in-hand view are individually back-projected and transformed into the base frame using the camera extrinsics obtained from hand-eye calibrations using EasyHeC++~\cite{hong2024easyhec++}.
    
    \item \textbf{Multi-view merging}: The two transformed point clouds are concatenated into a single dense colored cloud containing $200\text{k}-400\text{k}$ points.
    
    \item \textbf{Workspace cropping}: A fixed workspace box removes irrelevant background: $x\in[-0.5, 0.82]$, $y\in[-0.7, 0.8]$.
    
    \item \textbf{Geometry-preserving downsampling}: The cropped cloud is reduced to \textbf{8,192 points} using \textbf{farthest point sampling (FPS)} during training, which preserves coverage of manipulable surfaces and object geometry.
    
    \item \textbf{Real-time inference}: During deployment, point clouds are generated online at $\sim$30\,Hz. CDM is disabled for latency; raw RGB--D frames are fused, transformed, merged, cropped, and \textbf{uniformly sampled to 8,192 points}.
\end{enumerate}

\begin{figure}[tb]
    \centering
    \includegraphics[width=0.8\linewidth]{figures/realtime-pointcloud.png}
    \caption{\textbf{Real-time fused point cloud used during inference.}  
    Point clouds are updated at $\sim$30\,Hz and follow the same merging, base-frame transformation, and cropping conventions used in the training pipeline.}
    \label{fig:realtime_pc}
\end{figure}

\paragraph{Point Cloud Visualization.}
We include representative visualizations to illustrate the point clouds used during training and inference.

\paragraph{Training Setup.}
We collect \textbf{50 human teleoperated demonstrations per task}. Raw depth maps are enhanced using the \textbf{Camera Depth Model (CDM)}~\cite{liu2025manipulation}. CDM-refined depth is fused with RGB into colored point clouds, transformed to the robot base frame, merged, workspace-cropped, and \textbf{subsampled to 8,192 points via FPS}. All models are trained for \textbf{1,000 epochs}.



\section{Pre-training on large scale 3d datasets}

\subsection{Pre-training data}

Our pre-training corpus is primarily anchored by indoor scene scans (ScanNet\cite{dai2017scannet} and ARKitScenes\cite{baruch2021arkitscenes}), while being augmented by a synthetic dataset with part-level annotations (PartNeXt\cite{wangpartnext}). To align with the input configuration of the proposed framework, we downsampled all point cloud samples to 1024 or 8192 points. Specifically, as representing comprehensive indoor scenes with only 1024 or 8192 points results in excessive sparsity that impedes effective training, we employed a block partitioning strategy. This involves slicing large-scale scenes into smaller blocks to preserve local geometric details. Furthermore, the inclusion of PartNeXt is motivated by its multi-granularity segmentation, which provides rich geometry of individual objects, thereby facilitating the model's transferability to robotic perception tasks.

\subsection{Pre-training details}

We followed the training protocol established by PointSAM\cite{zhou2024point} to pre-train our 3D visual encoder. While PointSAM originally employs Voronoi tessellation for point cloud patchification, we empirically found that a k-Nearest Neighbor (k-NN) based grouping strategy aligns better with our specific data distribution and objectives; thus, we replaced the Voronoi module with k-NN to enhance local feature aggregation. Furthermore, the original PointSAM leverages the Uni3D~\cite{zhou2023uni3d}-large Vision Transformer (ViT-large) backbone to process high-resolution point clouds. In contrast, since our framework normalizes inputs to a fixed density of 1024 or 8192 points, utilizing such a heavy backbone would introduce unnecessary computational redundancy. Consequently, we adopted the more lightweight ViT-tiny, ViT-small, and ViT-base variants as the ViT backbones. Correspondingly, we scaled down the embedding dimensions and related hyperparameters to match this lightweight architecture. The PointSAM model was trained for 40 epochs on 8 NVIDIA H20 GPUs. We then utilized the pre-training weights to initialize our encoders.

\section{Implementation Details in Simulation Experiments}
\subsection{RoboTwin 2.0~\cite{chen2025robotwin} ``Easy''}
\begin{itemize}
    \item \textbf{RDT}~\cite{liu2025rdt1bdiffusionfoundationmodel} is pretrained for 100,000 steps with a batch size of 16 per GPU on 8 GPUs, and all single-task fine-tuning was conducted for 10,000 steps with a batch size of 16 per GPU on 4 GPUs.
    
    \item \textbf{Pi0}~\cite{black2024pi_0} is pretrained for 100,000 steps with a batch size of 32, and all fine-tuning was performed for 30,000 steps using the same batch size.

    
    \item \textbf{DP}~\cite{chi2024diffusionpolicyvisuomotorpolicy} is trained for 600 epochs with a batch size of 128 and a planning horizon of 8.
    
    \item \textbf{DP3}~\cite{ze20243d} is trained for 3,000 epochs with a batch size of 256, using a planning horizon of 8 and a point cloud resolution of 1,024, with precise segmentation of the background and tabletop.
    \item \textbf{Spatial Forcing}~\cite{li2025spatialforcingimplicitspatial} employs $\pi_0$ on Robotwin 2.0~\cite{chen2025robotwin} as the base model, which uses the PaliGemma as the VLM backbone. We train Spatial Forcing based on $\pi_0$ with LoRA on 1 NVIDIA H100 for 30k iterations. To ensure fairness, all training and evaluations follow the official settings.
    \item \textbf{DP+eVGGT} utilizes the VGGTS encoder (initialized with distilled weights from RoboTwin 2.0~\cite{chen2025robotwin}) as the frozen vision backbone within the UNet-based Diffusion Policy. The model is configured with a horizon of 8, 3 observation steps, and 6 action steps, using down-sampling dimensions of [256, 512, 1024]. It was trained for 100 epochs with a batch size of 48.
    \item \textbf{Ours} is trained for 1000 epochs with a batch size of 256, using a planning horizon of 16 and a point cloud resolution of 1024.
\end{itemize}
\subsection{RoboTwin 2.0~\cite{chen2025robotwin} ``Hard''}
\begin{itemize}
    \item \textbf{DP+evggt}~\cite{vuong2025improvingroboticmanipulationefficient} utilizes the VGGTS encoder (initialized with distilled weights from RoboTwin 2.0~\cite{chen2025robotwin}) as the frozen vision backbone within the UNet-based Diffusion Policy. The model is configured with a horizon of 8, 3 observation steps, and 6 action steps, using down-sampling dimensions of [256, 512, 1024]. It was trained for 60-80 epochs with a batch size of 48.
    \item \textbf{ACT+evggt}~\cite{vuong2025improvingroboticmanipulationefficient} is trained under a unified setup with a chunk size of 50, batch size of 8, and single-GPU training for 6,000 epochs. During deployment, we applied temporal\_agg for temporal aggregation to improve execution stability.
    \item \textbf{DP3}~\cite{ze20243d} is trained for 3,000 epochs with a batch size of 256, using a planning horizon of 8 and a point cloud resolution of 1,024, with precise segmentation of the background and tabletop.
    \item \textbf{RDT}~\cite{liu2025rdt1bdiffusionfoundationmodel} is pretrained for 100,000 steps with a batch size of 16 per GPU on 8 GPUs, and all single-task fine-tuning was conducted for 10,000 steps with a batch size of 16 per GPU.
    \item \textbf{ManiFlow}~\cite{yan2025maniflowgeneralrobotmanipulation} adopts a Transformer-based architecture (DiTX block) with 12 layers, 8 attention heads, and an embedding dimension of 768. It utilizes a DP3Encoder (PointNet-based) to process point cloud inputs with 128 downsampled points. The model was trained for 1,000 epochs with a batch size of 384 on a single 24GB NVIDIA RTX 4090 GPU. 
    \item \textbf{Ours} is trained for 1000 epochs with a batch size of 256, using a planning horizon of 16 and a point cloud resolution of 1024.
\end{itemize}
\subsection{Maniskill2~\cite{gu2023maniskill2}}
\begin{itemize}
\item \textbf{DP+SPuNet}~\cite{zhu2024point} employs the SPuNet34 architecture (adapted from Pointcept) as the encoder within the official UNet-based Diffusion Policy framework. It is trained for 1,800 epochs with a batch size of 32, a prediction horizon of 16, 8 action steps, and 2 observation steps. The diffusion network uses down-sampling dimensions of [512, 1024, 2048] and a diffusion step embedding dimension of 128.

\item \textbf{DP+PointNet}~\cite{zhu2024point} utilizes the standard PointNet encoder integrated with the Diffusion Policy. Adhering to the unified diffusion settings, it is trained for 1,800 epochs using the AdamW optimizer (weight decay 0.05) and OneCycle scheduler. The policy operates with a horizon of 16, performed on a single NVIDIA RTX 3090/4090 GPU.

\item \textbf{ACT+PonderV2}~\cite{zhu2024point} combines the Action Chunking Transformer with the PonderV2 backbone (which adopts the SPuNet34 architecture). The ACT policy is configured with 4 encoder layers, 7 decoder layers, a hidden dimension of 512, and a latent dimension of 32 with a KL divergence weight of 10.0. It is trained for 500 epochs with a chunk size of 100 (for ManiSkill2~\cite{gu2023maniskill2}) and applies temporal ensembling with an exponential decay factor of $k=0.01$ during inference.

\item \textbf{DP3}~\cite{ze20243d} is trained for 1,000 epochs with a batch size of 2048 without point cloud rgb, using a planning horizon of 8 and a point cloud resolution of 1,024, with precise segmentation of the background.

\item \textbf{Ours} is trained for 1000 epochs with a batch size of 2048 with point cloud colors, using a planning horizon of 16 and a point cloud resolution of 1024.
\end{itemize}









% ---- Bibliography ----
%
% BibTeX users should specify bibliography style 'splncs04'.
% References will then be sorted and formatted in the correct style.
%
\bibliographystyle{splncs04}
\bibliography{main}
\end{document}
